\section{An optimal subtraction and an explicit elliptic consequence}

Set $m=\sigma-\tfrac12>\tfrac12$, $x=\beta-\tfrac12$, and let $t$ be the vertical distance from the evaluation point to a zero. The paired real logarithmic-derivative kernel is
\[
 K_m(x,t)=\frac{m-x}{(m-x)^2+t^2}
             +\frac{m+x}{(m+x)^2+t^2}.
\]
It accounts for a zero together with its horizontal reflection. The classical uniform subtraction with $\kappa=1/\sqrt5$ is due to Stechkin; see H.~Kadiri, \emph{Explicit zero-free regions for Dedekind zeta functions}, Lemma~2.1, arXiv:1106.1868. Here we optimize over positive mixtures at each fixed real abscissa.

\paragraph{Theorem.}
Let $N^2=2m^2+m+\tfrac14$, $N>0$, and $k=m/N$. Then
\[
 K_m(x,t)-kK_N(x,t)\ge0
 \quad (|x|\le\tfrac12,\ t\in\mathbb R).
\]
For $0\le x\le\tfrac12$, at $t=0$ the difference is at least $1/(m-x)$. Equality holds at $x=\tfrac12$. The total weight of every finite nonnegative mixture $\sum_i w_iK_{n_i}$, with $n_i>\tfrac12$, satisfying these two requirements is at most $k$. In particular, $k>1/\sqrt5$ for $\sigma>1$.

\paragraph{Proof.}
For $X=m^2-x^2+t^2$ and $C=4x^2t^2$, one has $K_m=2mX/(X^2+C)$. Put $Y=N^2-x^2+t^2$. Since $Y\ge X>0$, the kernel comparison reduces to $XY\ge C$. Indeed
\[
 XY-C=(m^2-x^2)(N^2-x^2)+(m^2+N^2-6x^2)t^2+t^4\ge0.
\]
Here $m^2+N^2\ge3/2$ and $x^2\le1/4$. At zero distance the stronger inequality reduces to $2m(m+x)\le N^2-x^2$. The margin equals $(1/2-x)(2m+1/2+x)$, which vanishes at $x=1/2$. Finally $5m^2-N^2=(m-1/2)(3m+1/2)>0$.

For the converse, at $x=0$ multiply the first inequality by $t^2/2$ and let $|t|\to\infty$. This gives $\sum_i w_i n_i\le m$. At $x=1/2,t=0$ the stronger inequality gives
\[
 \sum_i w_i f(n_i)\le \frac1{m+1/2},\qquad
 f(v)=\frac1{v-1/2}+\frac1{v+1/2}.
\]
The graph of $f$ lies above its tangent at $N$, by convexity of $1/v$ on $v>0$. Taking $A=-f'(N)>0$, multiplying $f(N)+AN\le f(n_i)+An_i$ by $w_i$, and summing gives
\[
 (f(N)+AN)\sum_iw_i\le(m+1/2)^{-1}+Am=k(f(N)+AN).
\]
In the last step we used the exact boundary identity $kf(N)=(m+1/2)^{-1}$. This also proves that the unshifted zeta pole and its companion contribute exactly their singular charge $P/(\sigma-1)$ after subtraction. The proofs in \texttt{Stechkin.lean} implement the limiting argument by a finite error bound, so they require no interchange of an infinite sum and a limit.

\paragraph{Corollary for the original conductor-eleven function.}
Write $L=\log|T|$. Under the explicit-formula and elliptic local-factor inputs of the conductor-eleven paper, every nontrivial unitary zero $\beta+iT$ of $11.\mathrm{a1}$ with $|T|\ge262144$ satisfies
\[
 \min(\beta,1-\beta)>
 \frac1{13.518 L}+\frac1{25L^2}>
 \frac1{13.518 L-7}.
\]

\paragraph{Proof.}
In the archived five-frequency square take $P=413/900$, $g=9161/3600$, $b=0.15535$, $u=b/L$, $m=1/2+u$. Here $u\le1/50$. The square and Stechkin pair argument give
\[
 \frac1{u+d}\le P/u+(1-k)gL+(1-k)(M-P)+0.004,
 \qquad d=1-\beta,
\]
on the right of the critical line. The optimal abscissa $N+1/2$ is between $\sigma$ and $2$. The archimedean charge $M=C-g\log(2\pi)$ uses the same reference constant $C$ as the previous paper, in which $+P$ paid for the regular static pole. The exact companion cancellation removes this $+P$. The previous archimedean estimates apply unchanged for both abscissae. The included outward-rounded evaluation gives $M<-0.714$ and bounds the combined shifted gamma and regular-pole errors by $0.000029$. All finite-place inequalities and the applicability of the Hadamard identity to these functions are analytic inputs.

We give details of the small gain, since it applies to the original function at each specified ordinate. For $0<u\le0.02$,
\[
 1/\sqrt5+u/3\le k\le0.46.
\]
Indeed $m\le0.52$, $N^2\le21/16$, and comparison of squares gives $k\le0.46$. The identity
$5N^2(k+1/\sqrt5)(k-1/\sqrt5)=2u+3u^2$, together with
$5N^2(k+1/\sqrt5)<6$, proves the lower bound. Thus $(1-k)(M-P)+0.004\le-0.629$.
Put $k_0=0.44721359<1/\sqrt5$ and $q=P+(1-k_0)gb$. After multiplication by $u(u+d)$, the preceding bound gives
\[
 u\le(q-y)(u+d),\qquad y=u(0.629+gb/3).
\]
In particular $q-y>0$. Since $(q-y)(1/q+y/q^2)=1-y^2/q^2\le1$,
\[
 d\ge\frac{b/q-b}{L}+
         \frac{b^2(0.629+gb/3)}{q^2 L^2}.
\]
The two rational coefficients exceed $1/13.518$ and $1/25$, respectively. For $L\ge62/5$, clearing positive denominators verifies the last displayed improvement of the corollary. Reflection gives the opposite edge. In \texttt{StechkinWidth.lean}, the explicit formula is the named hypothesis \texttt{OriginalMajorant}. The real algebra, including both last rational comparisons, is proved by the kernel. The gamma evaluation is replayed by \texttt{elliptic\_reference.py}. The result improves the finite-height width and leaves its leading coefficient unchanged.
